% Reconstructed text from V5 PDF page 46; non-native source
\par\smallskip\noindent\rule{\textwidth}{0.4pt}\par{\sffamily\footnotesize Original V5 page 46 -- recovered text}\par\begin{flushleft}\footnotesize\ttfamily
\par\smallskip
\par\smallskip
\noindent Diagnostics and Serving Infrastructure for MAS 46\par
\par\smallskip
\par\smallskip
\par\smallskip
\noindent the serving cost, implement selective retention that protects only high-value diagnostic blocks,\par
\noindent or redesign the trace storage to decouple D from I (converting to Case 1).\par
\par\smallskip
\noindent 8.5 Relationship to ALEMC\par
\noindent The DRICE framework is complementary to ALEMC (Accuracy, Latency, Energy, Memory,\par
\noindent Carbon) proposed for on-device RAG evaluation. Where ALEMC focuses on single-system\par
\noindent performance metrics with explicit sustainability dimensions, DRICE focuses on cross-system\par
\noindent characterization — what each system assumes about, produces for, and depends on from adja-\par
\noindent cent layers. The two frameworks can be composed: ALEMC dimensions apply within each\par
\noindent DRICE component (e.g., the C dimension’s compute cost can be further decomposed into\par
\noindent ALEMC’s latency and energy metrics). A system’s DRICE profile identifies which layers it\par
\noindent spans; ALEMC quantifies how well it performs within each layer.\par
\par\smallskip
\noindent 8.6 Limitations of DRICE\par
\noindent DRICE is a characterization framework, not a ranking system. It does not produce a single score\par
\noindent because the dimensions are incommensurable: a high D score (rich evidence consumption) is\par
\noindent not inherently better than a high R score (sophisticated retention) — the relative importance\par
\noindent depends on the deployment context. The framework assumes that each dimension can be inde-\par
\noindent pendently assessed, which may not hold when dimensions interact (e.g., a higher R score may\par
\noindent increase C costs). We present DRICE as a structured vocabulary for cross-layer comparison,\par
\noindent not as a quantitative evaluation methodology.\par
\noindent  Several extensions could strengthen the framework. First, the E dimension (evidence depth)\par
\noindent currently uses our three-tier coding scheme; a more granular scale that captures the distinc-\par
\noindent tion between verified numbers and verified mechanisms would improve comparative precision.\par
\noindent Second, the C dimension aggregates four cost components; future work could weight these\par
\noindent components by their binding-constraint relevance in specific deployment scenarios. Third, the\par
\noindent framework does not capture temporal dynamics: a system’s DRICE profile may change as\par
\noindent workload patterns shift, memory pressure varies, or the diagnostic horizon extends. A dynamic\par
\noindent DRICE assessment that tracks profile changes over time would better capture the reality of\par
\noindent production deployments.\par
\noindent  The framework’s primary value lies in making implicit assumptions explicit. When a re-\par
\noindent searcher proposes a new diagnostic method, DRICE asks: what infrastructure assumptions does\par
\noindent it make (I)? When an infrastructure researcher designs a new eviction policy, DRICE asks: does\par
\noindent it consider diagnostic value (D)? This structured questioning is more actionable than a generic\par
\noindent call for “cross-layer optimization.”\par
\par\smallskip
\par\smallskip
\noindent 9 Discussion\par
\noindent 9.1 Key Patterns\par
\noindent Five patterns recur across the diagnostic and infrastructure layers:\par
\par\smallskip
\noindent The replay-accuracy trade-off is disappearing. Among the 12 failure-localization methods,\par
\noindent only two require environment replays (FAMAS, k=20; DoVer, 4–12 per case). The remaining\par
\end{flushleft}
